\documentclass{swmcmthesis}
% 数维杯官方模板（2026 秋季赛）。与 apmcmthesis 的三处接口差异：
%   1) 摘要用 \abstract{...} 宏，不是 \begin{abstract}...\end{abstract}；
%   2) 选题用 \zubie{ABCD}，不是 \tihao{A}；
%   3) 页眉两行（队号一行、页码一行），review_bench.HEADER_RE 已同时认两种排版。
% 宏参数里不要出现空行（空行会被当 \par，参数没加 \long 时直接报错），换行用 \\。
\title{<<title>>}
\baominghao{SWBMCM2634892}
\zubie{C}
\abstract{
\textbf{Background and Task} <<one sentence with a number>>. This paper addresses <<N>> problems: <<task 1>>; <<task 2>>; <<task 3>>. <<N>> models are developed. \\
\textbf{Modeling Framework} <<Q1 model>>, based on which <<Q2 model>>, and extending to <<Q3 model>>. \\
\textbf{Results per Question} <<Q1 number + conclusion>>; <<Q2 number + conclusion>>; <<Q3 number + conclusion>>. \\
\textbf{Validation and Limitation} <<validation 1 with number>>; <<validation 2 with number>>. The main limitation is that <<one concrete limitation with a number>>.
}
\keywords{<<model names in the order they appear above>>}
\date{2026}
\begin{document}
% 第 1 页 = 标题 + 摘要 + 关键词（\maketitle 生成），第 2 页 = 目录，第 3 页起正文。
\maketitle
\tableofcontents
\newpage
% 必须重置：本类默认不重置，正文会从物理第 3 页起编号（实测 "Page 3 of 5"），
% 于是 body_pages 判成 0、R-V02 报跳号。重置后与 apmcmthesis 同语义（Page 1 of N 指正文页）。
\setcounter{page}{1}

\section{Introduction}
\subsection{Background}
<<why this matters now, one paragraph>>
\subsection{Restatement of the Problems}
<<restate each question in your own variables>>
\subsection{Our Work}
<<what this paper delivers, one line per question>>
\section{Problem Analysis}
<<the mechanism, the key tension, what you simplify and why>>
\section{Assumptions and Justifications}
\begin{enumerate}
  \item <<assumption>> --- <<why>>; used by Section~<<x>>.
  \item <<assumption>> --- <<why>>; used by Section~<<x>>.
\end{enumerate}
\section{Notations}
\begin{table}[htbp]\centering\caption{Notations}\label{tab:notation}
\begin{tabular}{@{}lll@{}}\toprule Symbol & Meaning & Unit \\ \midrule
$<<sym>>$ & <<meaning>> & <<unit>> \\ \bottomrule\end{tabular}\end{table}
\section{Data Description and Preprocessing}
\subsection{Sources and Acquisition}
<<institution, indicator, period, URL, access date --- one line per source>>
\subsection{Cleaning}
<<missing rate, outliers, unit conversion, post-processing row count>>
\subsection{Rationality and Consistency}
<<cross-check two independent sources, report the gap in per cent>>
\section{Question 1: <<What it does>>}
\subsection{Data Insight}
<<fill>>
\subsection{Modeling and Solution}
<<fill>>
\subsection{Robustness and Conclusion}
<<fill>>
\section{Question 2: <<What it does>>}
\subsection{Data Insight}
<<fill>>
\subsection{Modeling and Solution}
<<fill>>
\subsection{Robustness and Conclusion}
<<fill>>
\section{Sensitivity Analysis and Error Analysis}
<<perturbation / resampling / variance decomposition; report the unfavourable half>>
\section{Conclusion}
<<per-question summary, each pointing at a number>>
\section{Strengths and Weaknesses}
\subsection{Strengths}
\begin{itemize}
  \item <<strength 1, points to a result>>
  \item <<strength 2>>
  \item <<strength 3>>
\end{itemize}
\subsection{Weaknesses}
\begin{itemize}
  \item <<weakness 1, with a number>>
  \item <<weakness 2, with a number>>
  \item <<weakness 3, with a number>>
\end{itemize}

\begin{thebibliography}{99}
\bibitem{ref1} <<author>>. <<title>>[J]. <<journal>>, <<year>>, <<vol>>(<<no>>): <<pages>>.
\bibitem{ref2} <<author>>. <<book>>[M]. <<place>>: <<publisher>>, <<year>>.
\bibitem{ref3} <<institution>>. <<report>>[R]. <<year>>.
\bibitem{ref4} <<institution>>. <<dataset>>[DB/OL]. <<url>>, accessed <<YYYY-MM-DD>>.
\end{thebibliography}

\appendix
% 必须用 \section{Appendix A: Code}：apmcmthesis.cls 里 \section 被改成罗马数字（I. ），
% \appendix 不会把它变成 "Appendix A"。而 gauntlet 认正文首页的锚点正是字面 "Appendix A"。
\section{Appendix A: Code}
% 先把 code/main.py 写出来，再把下一行取消注释。模板默认不启用：
% 缺文件时 listings 抛 "! Package Listings Error: File `code/main(.py)' not found"
% + "! Emergency stop."，exit=1，且 lastpage 跳号让页眉变成 "Page N of ??"。
%\lstinputlisting{code/main.py}
<<paste the full source of every script that produced a number, one appendix per script>>
\end{document}
